% ============================================================
\setpartpreamble{\partintro
  {Welche Bausteine in einen gemeinsamen Kern gehören, welche in Provider-Pakete, und wie ein
  zweiter Provider auf demselben Kern entsteht.}
  {
    \item Kapitel~\ref{ch:what-belongs-in-the-core}: Was in den Kern gehört
    \item Kapitel~\ref{ch:a-second-provider-sql-server}: Ein zweiter Provider: SQL Server
  }}
\part{Mehrere Provider}\label{part:providers}

% ------------------------------------------------------------
\chapter{Was in den Kern gehört}\label{ch:what-belongs-in-the-core}
\versionrisk{stable}

Ein zweiter Provider ist der Härtetest für die Paketarchitektur. Er zeigt, welche Teile
wirklich provider-unabhängig sind, und der erste Eindruck täuscht: Ein Baustein kann weder
SQL-Syntax noch Provider-Typen kennen und trotzdem eine relationale Datenbank brauchen, weil er
mit ADO.NET-Verbindungen oder mit Tabellen arbeitet. Das Beispiel hat deshalb drei Schichten,
jede in einem eigenen Paket:

\begin{datatable}{H{1.21}H{0.45}H{1.34}}
  \textbf{Baustein} & \textbf{Schicht} & \textbf{Begründung} \\
  \tablehead
  \code{Tenant\allowbreak{}Id}, Konverter, JSON-Reader/Writer & Kern & Reine .NET-Typen und
  EF-Value-Converter \\
  \tablerule
  \code{ITenant\allowbreak{}Owned}, \code{ITenant\allowbreak{}Provider}, Attribute & Kern &
  Für Entitäts-Assemblies ohne Provider-Abhängigkeit \\
  \tablerule
  \code{Is\allowbreak{}Tenant\allowbreak{}Property(\allowbreak{})}, Kandidatensuche & Kern &
  Nur Metadaten \\
  \tablerule
  Laufzeit-Annotation, \code{Get\allowbreak{}Tenant\allowbreak{}Property\allowbreak{}Name} &
  Kern & Von den provider-neutralen Interceptors gelesen \\
  \tablerule
  SaveChanges-Interceptor & Kern & Arbeitet nur mit dem Change Tracker \\
  \tablerule
  Schutz von Massen-Updates (Query-Expression-Interceptor) & Kern & Prüft nur den
  LINQ-Ausdrucksbaum \\
  \tablerule
  Evaluatable-Filter mit Attribut & Kern & Erkennt Funktionsklassen aller Pakete \\
  \tablerule
  Prüfung beim Laden (Materialization-Interceptor) & Kern & Liest nur Modell, Context und
  Werte der Zeile \\
  \tablerule
  Value Generation, Logging, \code{Add\allowbreak{}Acme\allowbreak{}Tenancy} & Kern & Ohne
  Datenbankbezug \\
  \tablerule
  \code{Tenant\allowbreak{}Session\allowbreak{}Interceptor} (abstrakt) & Relational &
  Arbeitet auf der \code{Db\allowbreak{}Connection}; nur der SQL-Befehl unterscheidet sich je
  Provider \\
  \tablerule
  \code{Tenant\allowbreak{}Table\allowbreak{}Mapping\allowbreak{}Validator} & Relational &
  Tabellen und Spalten: TPT, TPC, Entity Splitting \\
  \tablerule
  \code{Database.\allowbreak{}Apply\allowbreak{}Tenant\allowbreak{}Session(\allowbreak{})} &
  Relational & Nutzt den Session-Interceptor des Providers \\
  \tablerule
  Type Mapping, Translators, Migrations, RLS-Conventions & Provider & Store-Typen,
  SQL-Syntax, Schema-Objekte \\
  \tablerule
  Eigene SQL-Knoten, Aggregate & Provider & Abhängig von Syntax und Provider-Factory \\
\end{datatable}

Die Faustregel hat deshalb drei Stufen: \textbf{Was keine Datenbank braucht, gehört in den
Kern. Was eine relationale Datenbank, aber keinen SQL-Dialekt braucht, gehört in die
relationale Schicht. SQL-Syntax und Store-Typen gehören in den Provider.} Das Kernpaket
referenziert nur \code{Microsoft.\allowbreak{}Entity\allowbreak{}Framework\allowbreak{}Core},
das Paket
\code{Acme.\allowbreak{}Entity\allowbreak{}Framework\allowbreak{}Core.\allowbreak{}Tenancy.\allowbreak{}Relational}
ergänzt \code{Microsoft.\allowbreak{}Entity\allowbreak{}Framework\allowbreak{}Core.\allowbreak{}Relational},
und beide Provider-Pakete bauen darauf auf. \code{Package\allowbreak{}Layer\allowbreak{}Tests}
stellt sicher, dass kein relationaler Typ in den Kern zurückwandert.

\begin{note}[Die relationale Schicht von Anfang an trennen]
  Wer später einen nicht-relationalen Provider wie Cosmos DB unterstützen könnte, legt die
  relationale Schicht von Anfang an in ein eigenes Paket. Typen nachträglich in eine andere
  Assembly zu verschieben bricht die Binärkompatibilität, und Entitäts-Assemblies, die nur
  den Kern referenziert haben, verlieren die transitive Abhängigkeit auf
  \code{Microsoft.\allowbreak{}Entity\allowbreak{}Framework\allowbreak{}Core.\allowbreak{}Relational},
  auf die sie sich vielleicht für \code{To\allowbreak{}Table} oder
  \code{Has\allowbreak{}Column\allowbreak{}Type} verlassen haben. Das Beispiel hat genau diesen
  Schritt in Ausgabe~1.2 gemacht.
\end{note}

\begin{important}[Namenskonflikte zwischen Provider-Paketen]
  Zwei Provider-Pakete, die beide eine Extension-Method
  \code{Use\allowbreak{}Acme\allowbreak{}Tenancy} auf
  \code{Db\allowbreak{}Context\allowbreak{}Options\allowbreak{}Builder} definieren, sind
  nicht gleichzeitig referenzierbar: Der Aufruf wäre mehrdeutig. Das PostgreSQL-Paket nutzt
  den Builder von EF, das SQL-Server-Paket den Builder des Providers
  (\code{Use\allowbreak{}Sql\allowbreak{}Server(\allowbreak{}cs,\allowbreak{} x =>
  x.\allowbreak{}Use\allowbreak{}Acme\allowbreak{}Tenancy(\allowbreak{}))}). Wer mehrere
  Provider-Pakete plant, sollte von Anfang an konsequent den Provider-Builder verwenden, wie
  das NetTopologySuite-Plugin.
\end{important}

% ------------------------------------------------------------
\chapter{Ein zweiter Provider: SQL Server}\label{ch:a-second-provider-sql-server}
\versionrisk{medium}

Das Paket
\code{Acme.\allowbreak{}Entity\allowbreak{}Framework\allowbreak{}Core.\allowbreak{}Tenancy.\allowbreak{}Sql\allowbreak{}Server}
setzt dieselbe Idee mit den Mitteln von SQL Server um. Die Unterschiede zeigen, was ein
Plugin pro Datenbank tatsächlich neu lösen muss:

\begin{datatable}{H{0.81}H{1.00}H{1.19}}
  \textbf{Aspekt} & \textbf{PostgreSQL} & \textbf{SQL Server} \\
  \tablehead
  Mandant in der Session & \code{set\_\allowbreak{}config(\allowbreak{}name,\allowbreak{}
  wert,\allowbreak{} false)} &
  \code{sp\_\allowbreak{}set\_\allowbreak{}session\_\allowbreak{}context} \\
  \tablerule
  \code{Current\allowbreak{}Tenant(\allowbreak{})} in SQL &
  \code{nullif(\allowbreak{}current\_\allowbreak{}setting(\allowbreak{}.\allowbreak{}.\allowbreak{}.\allowbreak{}),\allowbreak{}
  '')::uuid} & \code{CAST(\allowbreak{}SESSION\_\allowbreak{}CONTEXT(\allowbreak{}N'…') AS
  uniqueidentifier)} \\
  \tablerule
  Spaltentyp & \code{uuid} & \code{uniqueidentifier} \\
  \tablerule
  RLS-Mechanismus & \code{CREATE POLICY} pro Tabelle & Predicate-Funktion + \code{CREATE
  SECURITY POLICY} \\
  \tablerule
  Lesen / Schreiben & \code{USING} / \code{WITH CHECK} & \code{FILTER} / \code{BLOCK
  .\allowbreak{}.\allowbreak{}.\allowbreak{} AFTER INSERT/\allowbreak{}UPDATE} \\
  \tablerule
  Tabelleneigentümer & Nur mit \code{FORCE} gebunden & Kein FORCE; Prädikate sind Teil der
  Abfrage \\
  \tablerule
  \code{DROP TABLE}, Umbenennen & Policies verschwinden mit bzw. folgen & Policy muss vorher weg und danach neu angelegt werden (SCHEMABINDING) \\
  \tablerule
  DDL in Batches & Transaktional, beliebig kombinierbar & \code{CREATE
  FUNCTION/\allowbreak{}SCHEMA} allein im Batch:
  \code{EXEC(\allowbreak{}N'.\allowbreak{}.\allowbreak{}.\allowbreak{}')} \\
  \tablerule
  Tabellenabbildungen & \multicolumn{2}{>{\RaggedRight\arraybackslash}p{0.6\linewidth}}{Gleicher
  Vertrag aus dem relationalen Paket (Kapitel~\ref{ch:conventions}, Seite~\pageref{sec:supported-table-mappings})} \\
  \tablerule
  Bewusste Ausnahme & Attribut oder
  \code{Has\allowbreak{}No\allowbreak{}Row\allowbreak{}Level\allowbreak{}Security(\allowbreak{})}
  & Nur das Attribut \\
\end{datatable}

\section{Aktivierung über den Provider-Builder}

\begin{sourcecode}[Anwendung und Infrastructure/\allowbreak{}Sql\allowbreak{}Server\allowbreak{}Tenancy\allowbreak{}Options\allowbreak{}Builder\allowbreak{}Extensions.\allowbreak{}cs (\allowbreak{}gekürzt, vereinfacht)]{csharp}
services.AddAcmeTenancy<HttpTenantProvider>();                    // Kernpaket
services.AddDbContext<ShopContext>((sp, o) => o
    .UseSqlServer(connectionString, sql => sql.UseAcmeTenancy(sp)));

public static SqlServerDbContextOptionsBuilder UseAcmeTenancy(
    this SqlServerDbContextOptionsBuilder sqlServerOptionsBuilder,
    string sessionKey = "TenantId",
    string policySchema = "acme")
{
    // Vom Provider-Builder zum eigentlichen DbContextOptionsBuilder
    var optionsBuilder =
        ((IRelationalDbContextOptionsBuilderInfrastructure)sqlServerOptionsBuilder)
        .OptionsBuilder;

    var extension = optionsBuilder.Options.FindExtension<SqlServerTenancyOptionsExtension>()
                    ?? new SqlServerTenancyOptionsExtension();
    // ... Settings validieren und setzen ...
    ((IDbContextOptionsBuilderInfrastructure)optionsBuilder)
        .AddOrUpdateExtension(extension.WithSettings(settings));

    return sqlServerOptionsBuilder;
}
\end{sourcecode}

\section{Migrations: Security Policy statt Policy-Ausdruck}

Die Convention vermerkt für SQL Server nur die Mandantenspalte; das SQL entsteht im
Generator. Die Predicate-Funktion ist für alle Tabellen dieselbe und wird idempotent
angelegt:

\begin{sourcecode}[Migrations/\allowbreak{}Internal/\allowbreak{}Sql\allowbreak{}Server\allowbreak{}Tenancy\allowbreak{}Migrations\allowbreak{}Sql\allowbreak{}Generator.\allowbreak{}cs (\allowbreak{}Auszug, vereinfacht)]{csharp}
protected virtual void CreatePolicy(string table, string? schema, string column,
    MigrationCommandListBuilder builder)
{
    var function = Delimit(PredicateFunctionName, Settings.PolicySchema);
    var target = Delimit(table, schema);
    var predicate = $"{function}({Delimit(column)})";

    // CREATE SCHEMA/FUNCTION müssen allein im Batch stehen -> EXEC(N'...')
    Statement(builder, $"IF SCHEMA_ID({Literal(Settings.PolicySchema)}) IS NULL "
        + $"EXEC({Literal($"CREATE SCHEMA {Delimit(Settings.PolicySchema)}")})");

    var functionSql =
        $"CREATE FUNCTION {function}(@TenantId uniqueidentifier) "
        + "RETURNS TABLE WITH SCHEMABINDING AS RETURN SELECT 1 AS [Allowed] "
        + $"WHERE @TenantId = CAST(SESSION_CONTEXT({Literal(Settings.SessionKey)}) "
        + "AS uniqueidentifier)";
    Statement(builder, $"IF OBJECT_ID({Literal(function)}, N'IF') IS NULL "
        + $"EXEC({Literal(functionSql)})");

    Statement(builder,
        $"CREATE SECURITY POLICY {PolicyName(table, schema ?? "dbo")} "
        + $"ADD FILTER PREDICATE {predicate} ON {target}, "
        + $"ADD BLOCK PREDICATE {predicate} ON {target} AFTER INSERT, "
        + $"ADD BLOCK PREDICATE {predicate} ON {target} AFTER UPDATE "
        + "WITH (STATE = ON)");
}

protected override void Generate(DropTableOperation operation, IModel? model,
    MigrationCommandListBuilder builder, bool terminate = true)
{
    // SCHEMABINDING: zuerst die Policy, gefunden über die Tabelle, die sie bindet
    DropPolicy(operation.Name, TableSchema(operation.Schema, model), builder);
    base.Generate(operation, model, builder, terminate);
}
\end{sourcecode}

\code{SCHEMABINDING} akzeptiert nur zweiteilige Namen. Eine Tabelle ohne explizites Schema
qualifizieren die Aufrufer deshalb mit dem Default-Schema des Modells oder, falls keines
gesetzt ist, mit \code{dbo}; sonst lehnt SQL Server die Security Policy ab.

\begin{important}[Verschachteltes SQL]
  Das \code{EXEC(\allowbreak{}N'.\allowbreak{}.\allowbreak{}.\allowbreak{}')} enthält SQL in
  einem String-Literal, das wiederum Literale enthält. Die Verschachtelung entsteht
  ausschließlich über \code{Generate\allowbreak{}Sql\allowbreak{}Literal} des
  String-Mappings, das Anführungszeichen korrekt verdoppelt; Bezeichner laufen über
  \code{Delimit\allowbreak{}Identifier}. Session-Schlüssel und Schema werden zusätzlich beim
  Konfigurieren validiert.
\end{important}

\section{Schemaänderungen an geschützten Tabellen}\label{sec:sqlserver-schema-bound-changes}

\code{SCHEMABINDING} schützt die Predicate-Funktion, bindet aber auch die Tabelle. SQL Server
lehnt \code{sp\_\allowbreak{}rename} für Tabelle und Mandantenspalte, \code{ALTER SCHEMA
.\allowbreak{}.\allowbreak{}.\allowbreak{} TRANSFER}, \code{ALTER COLUMN} und \code{DROP COLUMN}
ab, solange eine Security Policy die Tabelle bindet: „object participates in enforced
dependencies“. Pro Operation lässt sich das nicht lösen. EF sortiert die Operationen einer
Migration selbst, und beim Umbenennen der Mandantenspalte kommt die
\code{Rename\allowbreak{}Column\allowbreak{}Operation} \emph{vor} der
\code{Alter\allowbreak{}Table\allowbreak{}Operation}, die die geänderte Annotation trägt. Eine
reine Tabellenumbenennung erzeugt gar keine \code{Alter\allowbreak{}Table\allowbreak{}Operation}.

Der Generator überschreibt deshalb den Einstieg, der die ganze Operationsliste bekommt. Er
sucht alle Tabellen, deren gebundene Objekte sich ändern, verfolgt ihre Umbenennungen und
umklammert die Liste. Vorne entfernt er die Policies der betroffenen Tabellen unter ihren alten
Namen. Hinten legt er sie mit den endgültigen Tabellen-, Schema- und Spaltennamen aus dem
Zielmodell neu an:

\begin{sourcecode}[Generiertes SQL beim Umbenennen der Mandantenspalte (\allowbreak{}gekürzt)]{sql}
DECLARE @policy nvarchar(400) = (SELECT TOP (1) ... FROM sys.security_policies AS p
    WHERE EXISTS (SELECT 1 FROM sys.security_predicates AS s WHERE s.object_id = p.object_id
        AND s.target_object_id = OBJECT_ID(N'[dbo].[Invoice]')
        AND s.predicate_definition LIKE N'%[[]fn[_]tenant[_]predicate]%'));
IF @policy IS NOT NULL EXEC(N'DROP SECURITY POLICY ' + @policy);
EXEC sp_rename N'[Invoice].[TenantId]', N'Mandant', 'COLUMN';
EXEC sp_rename N'[Invoice].[IX_Invoice_TenantId]', N'IX_Invoice_Mandant', 'INDEX';
IF SCHEMA_ID(N'acme') IS NULL EXEC(N'CREATE SCHEMA [acme]');
IF OBJECT_ID(N'[acme].[fn_tenant_predicate]', N'IF') IS NULL EXEC(N'CREATE FUNCTION ...');
CREATE SECURITY POLICY [acme].[tenant_dbo_Invoice]
    ADD FILTER PREDICATE [acme].[fn_tenant_predicate]([Mandant]) ON [dbo].[Invoice], ...
\end{sourcecode}

Drei Entscheidungen stecken darin:

\begin{itemize}
  \item \textbf{Transaktion statt Ausschalten.} EF führt die Befehle einer Migration in einer
    Transaktion aus, und DDL ist in SQL Server transaktional. Scheitert ein Schritt, rollt
    SQL Server auch das Entfernen der Policy zurück. Eine dauerhaft entfernte oder mit
    \code{STATE = OFF} abgeschaltete Policy kann so nicht entstehen. Der Down-Weg ist dieselbe
    Umklammerung mit vertauschten Namen.
  \item \textbf{Entfernen über die Bindung, nicht über den Namen.} Gelöscht wird nur eine
    Policy, deren Prädikat die Predicate-Funktion des Plugins verwendet und die genau diese
    Tabelle bindet. Fremde Policies und Tabellen ohne RLS bleiben unberührt, und Policies
    unter einem früheren Namen werden trotzdem gefunden.
  \item \textbf{Schemaqualifizierte Namen.} Bis 1.0 hieß die Policy
    \code{tenant\_\allowbreak{}Tabelle}, sodass gleichnamige Tabellen in zwei Schemas
    kollidierten. Neue Policies heißen \code{tenant\_\allowbreak{}Schema\_\allowbreak{}Tabelle}.
    Bestehende Installationen behalten den alten Namen, bis die Tabelle das nächste Mal
    angefasst wird; dann ersetzt der Generator die Policy. Die verbleibende Mehrdeutigkeit
    (Schema \code{a\_b} mit Tabelle \code{c} gegen Schema \code{a} mit Tabelle
    \code{b\_c}) und Namen über 128 Zeichen lehnt die Convention beim Modellbau ab.
\end{itemize}

Grenzen: Das Wiederherstellen braucht das Zielmodell. Das ist der Normalfall, weil
\code{dotnet ef migrations add} zu jeder Migration eine Designer-Datei erzeugt. Fehlt es,
oder kommt eine geänderte Tabelle darin nicht vor, bricht der Generator ab, statt eine Policy
fallen zu lassen. Eigene schemagebundene Objekte der Anwendung (Views mit
\code{SCHEMABINDING}, eigene Security Policies) behandelt das Plugin nicht; sie blockieren
dieselben Operationen weiterhin. Datenoperationen, die EF zwischen die Schemaänderungen
sortiert, laufen innerhalb der Transaktion ohne Policy, also mit den Rechten der
Migrationsrolle. \code{Sql\allowbreak{}Server\allowbreak{}Migration\allowbreak{}Tests} prüft
mit von EF erzeugten Diffs gegen einen echten SQL Server: Anlegen, Spalte und Tabelle
umbenennen, Schemawechsel, jeden Schritt zurück und Löschen, nach jedem Schritt Lese- und
Schreibisolation über den App-Login. Dazu kommen gleichnamige Tabellen in zwei Schemas, eine
Policy unter dem alten Namen und die fremde Policy einer Tabelle ohne RLS.

\section{Session-Schlüssel und Policy-Schema ändern}\label{sec:sqlserver-session-settings}

\code{session\allowbreak{}Key} und \code{policy\allowbreak{}Schema} sind Options, wirken aber
an zwei Stellen: zur Laufzeit im Session-Interceptor, der den Schlüssel setzt, und in der
Datenbank in der Predicate-Funktion, die ihn liest. Bis 1.0 las der Generator beide Werte aus
den \emph{aktuellen} Options und legte die Funktion nur an, wenn sie fehlte. Ein Wechsel von
\code{TenantId} auf \code{NewTenant} erzeugte deshalb keine Migration. Der Interceptor setzte
danach den neuen Schlüssel, die Funktion las weiter den alten, und kein Mandant sah mehr Zeilen.
Schlimmer noch: Eine alte Migration hätte beim erneuten Ausführen, etwa für eine neue
Testdatenbank, Objekte mit den neuen Werten erzeugt.

Seit 1.1 trennt das Paket die beiden Rollen. Die Convention schreibt Schlüssel und Schema als
Modell-Annotationen \code{Acme:\allowbreak{}Sql\allowbreak{}Server\allowbreak{}Rls:\allowbreak{}Session\allowbreak{}Key}
und \code{Acme:\allowbreak{}Sql\allowbreak{}Server\allowbreak{}Rls:\allowbreak{}Policy\allowbreak{}Schema}
ins Modell und damit in den Snapshot. Der Annotation-Provider gibt sie an das relationale Modell
weiter, und der Differ erzeugt bei einer Änderung eine
\code{Alter\allowbreak{}Database\allowbreak{}Operation} mit altem und neuem Wert. Der Generator
liest beide Werte aus dem Zielmodell der Migration, die er gerade erzeugt; die Options gelten nur
noch für Modelle, die die Werte nicht festhalten. Damit bleibt jede Migration bei dem Stand,
mit dem sie erzeugt wurde.

\begin{datatable}{H{0.70}H{1.30}}
  \textbf{Fall} & \textbf{Verhalten der Migration} \\
  \tablehead
  Unverändert & Keine Operation, keine Phantomänderung \\
  \tablerule
  Schlüssel oder Schema geändert, alter Wert im Snapshot & In einer Transaktion: Policies aller
  geschützten Tabellen entfernen, alte Funktion löschen, Funktion mit den neuen Werten anlegen,
  Policies neu anlegen. Der Down-Weg macht dasselbe rückwärts. \\
  \tablerule
  Erste Migration eines Contexts oder erste nach dem Upgrade von 1.0 (alter Wert unbekannt) &
  Nichts wird gelöscht. Eine vorhandene Funktion muss den Schlüssel lesen, den das Modell nennt,
  sonst bricht die Migration mit \code{THROW 50000} ab. \\
  \tablerule
  Down auf einen Snapshot vor 1.1 & Funktion und Policies bleiben, wie sie sind. \\
\end{datatable}

Die dritte Zeile ist eine bewusste Grenze. Ohne festgehaltenen alten Wert lässt sich nicht
beweisen, dass eine vorhandene Funktion dem Plugin gehört und nicht einem anderen Context.
Eine bestehende Installation wechselt den Schlüssel deshalb in zwei Schritten: erst die erste
Migration nach dem Upgrade mit unverändertem Schlüssel, dann der Wechsel in einer eigenen
Migration.

Dieselbe Prüfung schützt mehrere Contexts in einer Datenbank. Die Funktion ist pro
Policy-Schema eine, und jede Policy des Plugins prüft nach dem Anlegen der Funktion deren
Definition. Konfiguriert ein zweiter Context einen anderen Schlüssel im selben Schema, scheitert
seine Migration, statt seine Tabellen an eine Funktion zu binden, die den falschen Schlüssel
liest. Contexts mit verschiedenen Schlüsseln brauchen verschiedene Policy-Schemas.

Laufzeit und Datenbank können trotzdem kurz auseinanderlaufen. Läuft eine neue App-Version mit
neuem Schlüssel, bevor die Migration eingespielt ist, setzt der Interceptor einen Schlüssel, den
die Funktion nicht liest. Die Folge ist kein Datenleck, sondern \emph{fail closed}: keine Zeilen,
und Schreibversuche scheitern am BLOCK-Prädikat. Die Migration gehört deshalb vor die neue
App-Version. \code{Sql\allowbreak{}Server\allowbreak{}Session\allowbreak{}Settings\allowbreak{}Tests}
prüft ohne Datenbank die Annotationen, den Snapshot, den leeren Diff, die erkannte Änderung und
die historische Migration. Gegen einen echten SQL Server prüft es den Wechsel von Schlüssel und
Schema mit Up und Down, das Upgrade einer 1.0-Installation und den zweiten Context.

\section{Laufzeit: derselbe Ablauf, anderer Befehl}

% source: src/Acme.EntityFrameworkCore.Tenancy.SqlServer/Runtime/SqlServerTenantSessionInterceptor.cs
\begin{sourcecode}[Runtime/\allowbreak{}Sql\allowbreak{}Server\allowbreak{}Tenant\allowbreak{}Session\allowbreak{}Interceptor.\allowbreak{}cs]{csharp}
public class SqlServerTenantSessionInterceptor : TenantSessionInterceptor
{
    public SqlServerTenantSessionInterceptor(ITenantProvider tenants)
        : base(tenants)
    {
    }

    protected override DbCommand CreateSessionCommand(
        DbConnection connection, DbContext? context, TenantId? tenant)
    {
        var key = context?.GetService<ISqlServerTenancySingletonOptions>()
                      .Settings.SessionKey
                  ?? SqlServerTenancySettings.Default.SessionKey;

        var command = connection.CreateCommand();
        command.CommandText = "EXEC sp_set_session_context @key = @key, @value = @value";
        AddParameter(command, "@key", key);
        AddParameter(command, "@value", tenant?.Value);    // NULL: keine Zeilen sichtbar
        return command;
    }
}
\end{sourcecode}

\section{Bewusst nicht portiert}

Das SQL-Server-Paket ist ein eigenständiger, aber schlankerer Provider. Fluent-API-Policies,
der optionale Query-Filter, eigene SQL-Knoten, Aggregate, Snapshot-Code-Generator und
Scaffolding fehlen dort. Das ist eine Produktentscheidung, kein technisches Hindernis: Jedes
dieser Kapitel lässt sich mit denselben Mustern auf SQL Server übertragen. Die Liste gehört
in die Dokumentation des Pakets, damit Nutzer wissen, was sie erwarten können.

% source: tests/Acme.EntityFrameworkCore.Tenancy.SqlServer.Tests/SqlServerTests.cs
\begin{sourcecode}[tests/\allowbreak{}.\allowbreak{}.\allowbreak{}.\allowbreak{}Sql\allowbreak{}Server.\allowbreak{}Tests/\allowbreak{}Sql\allowbreak{}Server\allowbreak{}Tests.\allowbreak{}cs (\allowbreak{}Auszug)]{csharp}
[Fact]
public async Task Block_predicate_rejects_rows_of_other_tenant()
{
    await using var db = ContextFor(TenantId.New());

    var ex = await Assert.ThrowsAsync<SqlException>(() => db.Database.ExecuteSqlAsync(
        $"INSERT INTO [Orders] ([TenantId], [Total]) VALUES ({Guid.NewGuid()}, 1)"));

    Assert.Equal(33504, ex.Number);      // Verstoß gegen BLOCK PREDICATE
}
\end{sourcecode}

Die Tests laufen wie bei PostgreSQL in zwei Ebenen: ohne Datenbank für Modell, Abfrage- und
Migrations-SQL und mit Testcontainers gegen einen echten SQL Server. Dort verbindet sich die
Anwendung mit einem eigenen Login statt als \code{sa}, um die Produktionssituation
nachzustellen.
